\section{Choix technologiques}
\label{sec:choix}

Nous travaillons en Python sur Google Colab, car aucun membre n'a de GPU, et le simulateur tourne en local. COCO 2017 est notre jeu de données proposé, parce qu'il est le seul candidat à couvrir les 6 classes avec des photos de rue réelles et des annotations de référence.

\subsection*{Jeu de données}

\begin{center}
\small
\begin{tabularx}{\textwidth}{>{\raggedright\arraybackslash}p{3.4cm}>{\raggedright\arraybackslash}p{2.2cm}YY}
  \toprule
  \textbf{Option} & \textbf{Classes couvertes} & \textbf{Atouts} & \textbf{Limites} \\
  \midrule
  \href{https://cocodataset.org/\#download}{COCO 2017}, recadrages par boîte (proposé) & Les 6 & Photos réelles, annotations de référence & Découpage à faire, fort déséquilibre attendu, beaucoup de petits objets \\
  CIFAR-10 + CIFAR-100 (plan B) & Les 6, en combinant les deux & Téléchargement simple, images déjà centrées & 32~×~32~px seulement, peu routier \\
  Pascal VOC 2012 & 5 sur 6 & Annotations propres, taille raisonnable & Pas de classe camion \\
  Jeux Kaggle \og véhicules \fg{} & Variable & Prêts à l'emploi & Qualité et licence inégales, souvent sans personne ni vélo \\
  \bottomrule
\end{tabularx}
\end{center}

Les chiffres réels par classe seront mesurés dans le notebook 01, avant toute décision sur les seuils.

\subsection*{Outils}

\begin{center}
\small
\begin{tabularx}{\textwidth}{>{\raggedright\arraybackslash}p{3.4cm}>{\raggedright\arraybackslash}p{4.6cm}Y}
  \toprule
  \textbf{Besoin} & \textbf{Outil retenu} & \textbf{Pourquoi} \\
  \midrule
  Calcul et notebooks & Google Colab, Jupyter, VS Code & GPU gratuit, même environnement pour les 3 membres \\
  Partage & Google Drive pour les données et les poids, GitHub pour le code & Un seul jeu de données ; pas de fichiers lourds dans Git \\
  Données et figures & NumPy, pandas, Matplotlib, seaborn & Exploration et graphiques du rapport \\
  Traitement d'images & OpenCV & Recadrage, redimensionnement et descripteurs HOG \\
  Modèles classiques & scikit-learn & SVM, Random Forest, k-NN et métriques d'évaluation \\
  CNN & PyTorch, torchvision & Architecture sur mesure, entraînée de zéro \\
  Carte du simulateur & OSMnx, OpenStreetMap & Vraies rues de Casablanca \\
  Affichage du simulateur & pygame, Panda3D & Vue 2D, et vue 3D sans GPU dédié \\
  Tests & pytest & Vérifier le simulateur à chaque modification \\
  \bottomrule
\end{tabularx}
\end{center}

L'approche classique sert de référence demandée par le sujet. Le CNN montre si l'apprentissage direct sur les images apporte un gain réel, malgré l'entraînement de zéro.
